% Deutsche Parallelfassung zu foundations/ontology.tex

\chapter{Begriffliche Struktur}
\label{chap:conceptual-structure}
\index{Begriffliche Struktur}

\section{Die Mehrdeutigkeit, die ein Datenmodell nicht stillschweigend lösen kann}

Nehmen wir an, die bearbeitete CAD-Kurve und eine Vermessungspolyline weichen
um mehrere Zentimeter voneinander ab. Ein Datenmodell kann beide
Koordinatensätze speichern, Identifier zuordnen und ein Residuum berechnen.
Aus diesen Operationen allein kann es nicht ableiten, was die Differenz
bedeutet. Die Vermessung kann das physische Gleis beobachten; die CAD-Kurve
kann einen beabsichtigten Zustand ausdrücken; beide können dasselbe Alignment
repräsentieren; und das Residuum kann nach dem anwendbaren Ingenieurwissen
akzeptabel sein oder nicht.

Jede Aufzeichnung als weitere Version „des Alignments“ zu behandeln, verbirgt
genau die Unterscheidungen, die zum Nachdenken über die Abweichung erforderlich
sind. Jede unterschiedliche Geometrie als anderes Objekt zu behandeln,
verliert Kontinuität. Die Messaufzeichnung als Realität zu behandeln,
verwechselt Evidenz mit ihrem Gegenstand. Ein berechnetes Bestehen oder
Durchfallen als Entscheidung zu behandeln, überträgt Autorität auf eine
Operation, die sie nicht besitzt.

Die minimale Abhilfe besteht darin, Objekt, Zustand, Realisierung,
Repräsentation, Beobachtung, Wissen, Bewertung und Entscheidung verschiedene
formale Rollen zuzuweisen. Nach ihrer Trennung können sie ausdrücklich
aufeinander bezogen werden: Eine Beobachtung kann einen physischen Zustand
betreffen; eine Repräsentation kann einen beabsichtigten Zustand offenlegen;
eine Bewertung kann beide unter anwendbarem Wissen vergleichen; eine
Entscheidung kann die resultierende Evidenz verwenden, ohne mit ihr identisch
zu werden.

\section{Das Engineering Object und seine Zustände}

Ein Engineering Object ermöglicht dauerhafte Unterscheidbarkeit über Wechsel
von Repräsentation, Realisierung, Arbeitsablauf und Projektbezug hinweg. Ein
Alignment ist daher nicht mit einer CAD-Kurve, einem GIS-Feature, einer
Austauschentität, einer Datenbankzeile oder einer gerenderten Linie identisch.
Mehrere Beschreibungen können dasselbe Alignment betreffen, und eine geänderte
Beschreibung entscheidet für sich genommen nicht, ob das Engineering Object
verändert oder ersetzt wurde.

Ein Zustand beschreibt dieses Objekt unter benannten ingenieurtechnischen
Qualifikationen. Beabsichtigte, metrische, physische, betriebliche und
bewertete Zustände beantworten unterschiedliche Fragen über ein Objekt.
Zustandsänderung kann daher erörtert werden, ohne Objektidentität zu einer
Nebenwirkung von Koordinatengleichheit zu machen.

Beim Alignment bilden intrinsische longitudinale Ordnung und
Krümmungsverlauf die konstruktive Grundlage, aus der ausgewählte geometrische
Zustände abgeleitet werden. Spätere Kapitel formalisieren diese Ableitungen;
die begriffliche Anforderung hier lautet, dass der abgeleitete Zustand nicht
zum Owner der Identität wird.

\section{Realisierung und Repräsentation}

Realisierung und Repräsentation beantworten unterschiedliche Fragen.

Metrische Realisierung fragt, wie eine intrinsische Konstruktion unter einem
ausdrücklichen metrischen Raum, einem Realization Context und anwendbaren
Regeln messbare räumliche Bedeutung erhält. Physische Realisierung fragt, wie
Ingenieurabsicht zu physischer Infrastruktur wird. Keine von beiden ist
lediglich ein Wechsel des Dateiformats.

Eine Repräsentation wählt Informationen aus, damit diese kommuniziert,
gespeichert, ausgetauscht, visualisiert oder verarbeitet werden können.
CAD-Geometrie, GIS-Datensätze, IFC-Entitäten, abgetastete Polylinien,
Diagramme und Laufzeitstrukturen sind Repräsentationen. Eine Repräsentation
kann einen realisierten Zustand offenlegen, aber keine fehlende
Realisierungssemantik erzeugen oder für sich allein Identität begründen.

Diese Trennung ermöglicht eine nützliche Frage: Entsteht eine Differenz aus
der beabsichtigten Konstruktion des Objekts, dem gewählten metrischen Kontext,
der physischen Realisierung, der Beobachtung oder lediglich aus der
Repräsentation?

\section{Beobachtung und Wissen}

Eine Beobachtung ist quellengebundene Information über ein Objekt, einen
Zustand oder Qualifikationen einschließlich der Provenienz, Methode, Zeit,
Bedingungen und Unsicherheit, die zu ihrer Interpretation als Evidenz
erforderlich sind. Eine Vermessungspolyline ist daher nicht einfach „das reale
Alignment“. Sie ist eine Beobachtung, deren Beziehung zu einem physischen
Zustand begründet werden muss.

Engineering Knowledge ist Information, die innerhalb eines ausdrücklichen
Geltungsbereichs und mit benannter Provenienz, Anwendbarkeit und Autorität zur
ingenieurtechnischen Nutzung angenommen wurde. Beobachtungen können Wissen
stützen; Speicherung oder Wiederholung nimmt sie nicht automatisch an. Normen,
Ingenieurprinzipien, abgeleitete Ergebnisse und verantwortete Entscheidungen
können Wissen ebenfalls stützen, ohne austauschbare Quellen zu werden.

Die Unterscheidung erlaubt, widersprüchliche Beobachtungen weiterhin
festzuhalten, während ein Ingenieur bestimmt, worauf für die gegenwärtige
Aufgabe vertraut werden darf.

\section{Bewertung und Entscheidung}

Bewertung wendet Wissen, Nebenbedingungen, Präferenzen und
Vergleichsverfahren auf Zustände oder Alternativen an. Sie kann ein Residuum
ableiten, eine Verletzung bestimmen, Kandidaten ordnen oder zeigen, wie sich
die Krümmungsänderung fortpflanzt. Diese Ergebnisse sind Evidenz innerhalb
eines Ingenieurprozesses.

Entscheidung ist die verantwortete Handlung, die innerhalb einer bestimmten
Autorität und unter benannten Qualifikationen annimmt, ablehnt oder auswählt.
Ein Solver, Viewer oder eine Rule Engine kann diese Handlung unterstützen,
übt sie aber nicht allein dadurch aus, dass er eine Ausgabe erzeugt. Diese
Grenze erlaubt leistungsfähigere Berechnung, ohne Verantwortung unklarer zu
machen.

\section{Verfügbare Beziehungen}

Nach der Trennung der Rollen kann die Thesis Abhängigkeiten ohne
Kategorienfehler formulieren:

\begin{itemize}
\item Ein Engineering Object kann mehrere Zustände und Repräsentationen
  besitzen;
\item intrinsische Konstruktion kann metrisch realisiert werden, ohne Identität
  auf Koordinaten zu übertragen;
\item ein physischer Zustand kann beobachtet werden, ohne dass die Beobachtung
  zu diesem Zustand wird;
\item eine Beobachtung kann durch ausdrückliche Interpretation und Annahme
  Wissen stützen;
\item eine Bewertung kann anwendbares Wissen verwenden und Evidenz für eine
  Entscheidung erzeugen, ohne zur Entscheidung zu werden.
\end{itemize}

Diese Beziehungen machen aus der früheren Abweichung eine behandelbare
Ingenieurfrage. Die beabsichtigten und beobachteten Aufzeichnungen können auf
ein Alignment bezogen, ihre Zustände und Provenienz bestimmt und ein Vergleich
unter benanntem Wissen und benannten Qualifikationen durchgeführt werden.

\section{Kanonische Interpretation}

\begin{proposition}
Engineering Objects, Zustände, Realisierungen, Repräsentationen,
Beobachtungen, Engineering Knowledge, Bewertungen und Entscheidungen besitzen
verschiedene Rollen. Ihre ausdrücklichen Beziehungen bewahren Kontinuität und
machen zugleich ausgewählte Folgen berechenbar.
\end{proposition}

Diese begriffliche Struktur erlaubt nun die formale Entwicklung. Die folgenden
Kapitel formulieren die Axiome, die Notation, die Zustandsmodelle und die
Operatoren, die erforderlich sind, um diese Beziehungen genau auszudrücken.
Sie dürfen eine Abhängigkeit formalisieren; sie dürfen die
Verantwortungsgrenze, die der Abhängigkeit Bedeutung gegeben hat, nicht
aufheben.
